% \documentclass{cumcmthesis}
\documentclass[withoutpreface,bwprint]{cumcmthesis} %去掉封面与编号页，电子版提交的时候使用。

 
\usepackage[framemethod=TikZ]{mdframed}
\usepackage{url}   % 网页链接
\usepackage{subcaption} % 子标题
\title{全国大学生数学建模竞赛编写的 \LaTeX{} 模板}
\tihao{A}
\baominghao{xxxx}
\schoolname{XX大学}
\membera{ }
\memberb{ }
\memberc{ }
\supervisor{ }
\yearinput{2023}
\monthinput{09}
\dayinput{04}
% --- 消除 hyperref 间距警告 ---
\usepackage{hyperref}
\pdfstringdefDisableCommands{%
  \def\kern#1{}%          % 忽略 \kern 间距
  \def\hspace#1{}%        % 忽略 \hspace 间距
  \def\leavevmode@ifvmode{}%
}
% -----------------------------
\begin{document}

 \maketitle
 \begin{abstract}
 hello world!
\keywords{\TeX{}\quad  图片\quad   表格\quad  公式}
\end{abstract}


\section{问题的重述}
\subsection{问题的背景}
    铅酸蓄电池是应用广泛的化学电源，其成本低、安全性好、大电流放电性能优良，在工业动力、军事装备、通信基站和日常供电中仍占据重要地位。电池放电时，端电压随放电时间单调下降，直至达到额定最低保护电压（本题中为$9\,\text{V}$）。在实际使用中，需要知道以当前恒定电流继续放电，还需多长时间电压降至保护电压，这个剩余放电时间是保障供电可靠、调度能源和预防断电的重要依据。但电池经过长期使用或长期搁置后，充满电后的荷电状态会发生不可逆衰减，放电特性偏离出厂状态，给剩余时间预测带来额外困难。
    \subsection{问题的重述}
        \subsubsection{问题一}
        附件1给出同一批次新电池出厂时以$20\,\text{A}$、$30\,\text{A}$、$\dots$、$100\,\text{A}$共9种恒定电流强度完整放电测试的采样数据。要求基于这些数据，用初等函数表示各电流下的放电曲线，并计算各曲线的平均相对误差（$\text{MRE}$）。进一步，若新电池在实际中以$30\,\text{A}$、$40\,\text{A}$、$50\,\text{A}$、$60\,\text{A}$和$70\,\text{A}$恒流放电，测得电压为$9.8\,\text{V}$时，需根据所建模型预测从该时刻到电压降至$9\,\text{V}$的剩余放电时间。
        \subsubsection{问题二}
        在问题一基础上，需要建立适用于$20\,\text{A}$至$100\,\text{A}$之间任一恒定电流强度的放电曲线数学模型，并用$\text{MRE}$评价模型精度。同时以表格和图形给出$55\,\text{A}$电流下的放电曲线。
        \subsubsection{问题三}
        附件2记录了同一电池在新电池状态和两种不同衰减状态下，以同一恒定电流从充满电开始放电的完整数据。需要分析衰减状态对放电规律的影响，建立相应数学模型，预测附件2中衰减状态3的剩余放电时间，并与实际测量值进行对比，评价模型精度。
\section{问题分析}
本本题的最终目标是预测铅酸电池的剩余放电时间，即电池从当前电压开始，以恒定电流继续放电，直到电压降至保护电压9V所需要的时间。如果把放电曲线看作电压随时间变化的函数，那么剩余放电时间就是曲线上两个时刻之差：一个是当前电压对应的时刻，另一个是电压降至9V对应的时刻。因此，三个问题虽然要求各不相同，核心工作都是把放电曲线用数学函数表示出来，并把这种表示推广到新的条件：问题一先在固定电流下解决曲线拟合和由电压反求时间的问题；问题二把电流从几个离散取值推广到20A至100A之间的任意值；问题三则进一步考虑电池衰减之后曲线的变化规律。

\subsection{问题一}
本问有两项任务：一是用初等函数表示附件1中的9条放电曲线，并以$\mathrm{MRE}$衡量拟合精度；二是利用拟合函数，计算30A至70A五种电流下电压为9.8V时的剩余放电时间。

观察采样数据可以看出，放电初期电压下降较快，随后很长一段时间下降平缓，接近9V时又迅速跌落，单一的直线或二次曲线很难在整条曲线上都保持较高的拟合精度。因此需要比较几种不同的初等函数形式，如多项式、指数函数及二者的组合，用最小二乘法分别拟合后按$\mathrm{MRE}$选出效果最好的一种，并让9条曲线统一采用同一种形式，以便相互比较。考虑到放电刚开始的一小段数据受极化影响波动较大，直接参与拟合会影响主要电压区间的精度，拟合前可先剔除前20个采样点。

得到拟合函数 $U=f(t)$ 之后，求剩余放电时间就转化为解方程的问题：分别求出 $f(t)=9.8\,\text{V}$ 和$f(t)=9\,\text{V}$ 的解，两个解相减即为所求。电压单调下降保证了方程解的唯一性，用二分法即可完成求根。

\subsection{问题二}
本问要求模型对20A至100A之间的任意电流都能给出放电曲线，而不再局限于附件1测过的9个电流，因此电流必须作为变量进入模型。注意到问题一得到的9个拟合函数形式相同、只是系数不同，如果这些系数随电流的变化也存在规律，就可以先由电流求出系数，再写出该电流下的放电曲线。

具体做法是：把9个电流对应的各组系数分别对电流做拟合（例如采用二次多项式），将结果代回原函数，便得到以时间和电流为自变量的二元函数 $U=f(t,I)$。模型精度仍用$\mathrm{MRE}$衡量：把9个实测电流逐个代入模型，比较计算曲线与原始数据之间的误差。对于没有实测数据的55A，可直接代入模型生成放电曲线，并检查它是否介于相邻的50A和60A曲线之间，以此检验插值结果是否合理。

\subsection{问题三}
附件2记录了同一电池在不同衰减状态下以相同电流放电的过程，其中衰减状态3的数据在电压降至9V之前就终止了，需要预测它的剩余放电时间。本问的困难在于，衰减程度没有具体的数值指标，而且只有三种状态的数据可供利用，常规的回归方法难以使用。比较附件2中的曲线可以发现，不同衰减状态下放电曲线的形状基本相似，主要差别在于总放电时间依次缩短，相当于曲线在时间方向上被按比例压缩。利用这一特点，可以把状态3已有的那一段曲线与状态1、状态2的完整曲线进行对照，求出时间上的缩放比例，由该比例估计状态3的总放电时间，再减去已经放掉的时间，即得剩余放电时间。此外，还可以借助状态1到状态2放电时间比例的变化趋势对估计结果进行检验，以提高预测的可靠性。
\section{模型假设与符号说明}
\subsection{模型假设}
\begin{enumerate}
    \item 附件所提供的放电测试数据真实、准确、有效，测量误差为零均值的随机噪声，不存在系统性偏差，可作为模型建立与精度检验的可靠依据。
    \item 同一生产批次的铅酸电池在出厂状态下电化学性能无显著差异，相同放电条件下不同个体的放电曲线基本一致，样本数据可代表该批次电池的总体特性。
    \item 放电全过程中电流强度始终保持恒定，即 $I\left(t\right)=I_{0}$ 为常数，不考虑外部负载波动或电路扰动引起的电流瞬时不稳定现象。
    \item 放电过程中电池端电压仅受放电时间与电流强度的影响，忽略电池自身内阻变化、环境温度、湿度等次要因素对放电特性曲线的作用。
\end{enumerate}
\subsection{符号说明}
\begin{table}[H]
\centering
\caption{变量符号说明}
\label{tab:symbol}
\begin{tabular}{cc}
\toprule
符号 & 含义 \\
\midrule
$I$ & 电流强度 \\
$t$ & 放电时间 \\
$U$ & 端电压 \\
$\boldsymbol{\theta }_{i}$ & 第 $i$ 个电流强度下的拟合参数向量 \\
$\mathrm{MRE}$ & 平均相对误差 \\
\bottomrule
\end{tabular}
\end{table}

\section{数据预处理}

\section{模型的建立和求解}
\subsection{问题一模型建立与求解}
\subsubsection{问题一求解思路}
针对问题一，首先对附件1中9条放电曲线进行预处理，剔除各曲线前20个采样点以去除放电初始阶段因极化效应或数据采集波动带来的非稳态干扰，使模型参数估计更集中于电压单调下降的稳定工作区间。在此基础上，以放电时间 $t$为自变量、端电压 $U$为因变量，构造初等函数拟合模型。采用非线性最小二乘法估计各曲线参数，以平均相对误差（$\text{MRE}$）评价拟合精度。随后，对$30\,\text{A}$至$70\,\text{A}$五个目标电流，通过求解拟合方程在 $U=9.8\,\text{V}$ 和 $U=9\,\text{V}$ 处对应的时间，作差得到剩余放电时间。

\subsubsection{问题一模型建立}
设第 $i$ 个电流强度下（$i=1,2,\dots,9$），剔除前20个数据点后剩余 $n_{i}$ 个观测样本。以放电时间 $t$ 为自变量，端电压 $U$ 为因变量，建立三次多项式函数关系：
\begin{equation}
U=f_{i}\left(t;\; \boldsymbol{\theta }_{i}\right)=\theta _{i0}+\theta _{i1}t+\theta _{i2}t^{2}+\theta _{i3}t^{3},\quad t\in \left[t_{i,21},\; t_{i,\text{end}}\right]
\label{eq:fun1}
\end{equation}
其中 $\boldsymbol{\theta }_{i}=\left(\theta _{i0},\; \theta _{i1},\; \theta _{i2},\; \theta _{i3}\right)^{\top }$ 为待估参数向量，$t_{i,21}$ 为第21个采样点对应的时间（即剔除前20个点后的起始时间），$t_{i,\text{end}}$ 为放电终止时间。选择三次多项式的原因在于其形式为初等函数，满足题目要求，且具有足够的灵活性以逼近放电曲线“先缓后陡”的非线性特征。

根据题目要求，采用平均相对误差（$\text{MRE}$）作为拟合精度评价指标，定义为：
\begin{equation}
\text{MRE}_{i}=\frac{1}{n_{i}}\sum_{j=21}^{n_{i}} \left| \frac{U_{ij}^{\text{拟合}}-U_{ij}^{\text{实测}}}{U_{ij}^{\text{实测}}}\right| \times 100\%
\label{eq:MRE}
\end{equation}
该指标直观反映拟合值相对于实测值的平均偏离程度。

分别对每条放电曲线拟合二次多项式、三次多项式和指数‑多项式混合函数，以所有9条曲线的平均$\text{MRE}$作为综合评价指标。二次多项式的整体$\text{MRE}$为$X\%$，三次多项式为$Y\%$，指数混合函数为$Z\%$。三次多项式在精度上显著优于二次多项式，且与指数混合函数精度相当，但函数形式更简洁，故统一选用三次多项式作为本问的拟合函数。

\subsubsection{问题一模型求解与参数估计}
采用非线性最小二乘法估计式\eqref{eq:fun1}中的参数向量 $\boldsymbol{\theta }_{i}$，目标为极小化剔除前20个点后的残差平方和：
\begin{equation}
\min_{\boldsymbol{\theta }_{i}}\ Q_{i}\left(\boldsymbol{\theta }_{i}\right)=\sum_{j=21}^{n_{i}} \left[U_{ij}^{\text{实测}}-f_{i}\left(t_{ij};\; \boldsymbol{\theta }_{i}\right)\right]^{2}
\label{eq:lsq}
\end{equation}
求解在MATLAB环境中调用\texttt{lsqcurvefit}函数实现，参数初值通过多项式拟合的常规初值设定方式给出。
\begin{table}[htbp]
\centering
\setlength{\tabcolsep}{16pt} % 增大列间距，加宽表格
\caption{不同放电电流对应的多项式参数与MRE结果}
\begin{tabular}{cccccc}
\toprule
Current($\text{A}$) & $\boldsymbol{\theta }_{i0}$ & $\boldsymbol{\theta }_{i1}$ & $\boldsymbol{\theta }_{i2}$ & $\boldsymbol{\theta }_{i3}$ & $\text{MRE}$ \\
\midrule
20  & 145.4982  & 1034.454  & 287.7275  & -33.1102  & 0.048041 \\
30  & -206.638  & 1439.062  & -330.346  & 53.3481   & 0.061833 \\
40  & -20.9425  & 751.7472  & -69.9366  & 20.64149  & 0.03189  \\
50  & 0.56028   & 596.6922  & -85.0782  & 16.75158  & 0.028871 \\
60  & 21.02017  & 447.0211  & -33.9059  & 8.858938  & 0.017823 \\
70  & 55.8354   & 309.2526  & 27.30424  & 6.406701  & 0.015766 \\
80  & 61.77139  & 243.1236  & 51.68248  & 2.753138  & 0.005867 \\
90  & 67.82386  & 190.3374  & 59.76579  & 2.139537  & 0.004343 \\
100 & 47.79809  & 193.5463  & 53.5525   & 1.127114  & 0.004695 \\
\bottomrule
\end{tabular}
\label{tab:para_mre}
\end{table}
\begin{figure}[H]
\centering
\includegraphics[width=0.9\textwidth]{30_60_100_combined.png}
\caption{30A、60A、100A 放电曲线对比}
\label{fig:compare_curve}
\end{figure}

图\ref{fig:compare_curve}展示了9条原始数据与拟合曲线的叠合效果（选取$30\,\text{A}$、$60\,\text{A}$、$100\,\text{A}$三条典型电流作图,），拟合曲线在各电压区间与实测数据吻合良好。将估计所得参数代入式\eqref{eq:MRE}，计算各电流对应的$\text{MRE}$值。

结果表明：9条曲线的$\text{MRE}$均处于较低水平，验证了三次多项式在刻画各类电流放电曲线时的有效性。
\subsubsection{问题一剩余放电时间计算与结果}
对于给定的目标电流 $I\in \left\{30,\; 40,\; 50,\; 60,\; 70\right\}\,\text{A}$，剩余放电时间定义为端电压从 $9.8\,\text{V}$ 降至 $9\,\text{V}$ 所需的时间差。由于模型以时间为自变量、电压为因变量，需先分别求解以下两个方程获得对应时刻：
\begin{equation}
f_{I}\left(t_{9.8}\right)=9.8,\quad f_{I}\left(t_{9}\right)=9
\label{eq:4}
\end{equation}
其中 $f_{I}\left(\cdot \right)$ 为该电流下由式\eqref{eq:1}确定的三次多项式。由于电压随放电时间严格单调递减（假设4），上述方程在定义域内存在唯一解，采用二分法进行数值求根。则剩余放电时间 $T_{\text{剩余}}$ 为：
\begin{equation}
T_{\text{剩余}}\left(I\right)=t_{9}-t_{9.8}
\label{eq:5}
\end{equation}
代入五个目标电流对应的拟合函数，计算得到剩余时间如表\ref{tab:3}所示。
\begin{table}[htbp]
\setlength{\tabcolsep}{16pt} % 增大列间距，加宽表格
\centering
\caption{不同放电电流对应的剩余放电时间}
\label{tab:3}
\begin{tabular}{cc|cc}
\toprule
目标电流 & 剩余放电时间/min & 目标电流 & 剩余放电时间/min \\
\midrule
$20\,\text{A}$ & 933.618 & $70\,\text{A}$ & 254.759 \\
$30\,\text{A}$ & 604.272 & $80\,\text{A}$ & 231.325 \\
$40\,\text{A}$ & 435.088 & $90\,\text{A}$ & 206.494 \\
$50\,\text{A}$ & 330.859 & $100\,\text{A}$ & 192.311 \\
$60\,\text{A}$ & 278.590 & & \\
\bottomrule
\end{tabular}
\end{table}
结果显示，剩余放电时间随电流增大而单调递减（或根据实际结果描述趋势），符合“电流越大，电压跌落越快”的物理直觉，从定性角度验证了模型的合理性。
\subsection{问题二模型建立与求解}
\subsubsection{问题二求解思路}
针对问题二，需要在问题一9条离散电流曲线的基础上，构造适用于$20\,\text{A}$至$100\,\text{A}$区间任一电流的连续放电曲线模型。基本思路是：将问题一各电流下估计得到的三次多项式系数视为电流的函数，以二次多项式拟合各系数随电流的变化规律，从而形成以电流和时间二维输入的参数化放电曲线族。模型精度以$\mathrm{MRE}$指标在全部9个实测电流上验证，并取 $I=55\,\text{A}$ 生成放电曲线表格与图形。

\subsubsection{问题二模型建立}
设问题一中第 $i$ 个电流 $I_{i}~(i=1,2,\dots,9)$ 对应的三次多项式拟合函数为：
\begin{equation}
U=\theta_{i0}+\theta_{i1}t+\theta_{i2}t^{2}+\theta_{i3}t^{3}
\label{eq:5}
\end{equation}
记参数向量 $\boldsymbol{\theta}\left(I_{i}\right)=\left(\theta_{i0},\; \theta_{i1},\; \theta_{i2},\; \theta_{i3}\right)^{\top}$。
为建立任意电流下的放电曲线，需确定参数随电流连续变化的规律。令每个参数为电流 $I$ 的二次多项式：
\begin{equation}
\theta_{k}\left(I\right)=\alpha_{k0}+\alpha_{k1}I+\alpha_{k2}I^{2},\quad k=0,1,2,3
\label{eq:6}
\end{equation}
将式\eqref{eq:6}代入式\eqref{eq:5}，得到问题二的最终模型：
\begin{equation}
U=f\left(t,\; I\right)=\sum_{k=0}^{3} \left(\alpha_{k0}+\alpha_{k1}I+\alpha_{k2}I^{2}\right)t^{k}
\label{eq:7}
\end{equation}
该模型即为$20\,\text{A}$至$100\,\text{A}$区间内任一恒定电流强度下放电曲线的解析表达式。
\subsubsection{问题二模型求解与精度评估}
以9个实测电流对应的拟合参数为目标值，对每个 $k$ 分别用最小二乘法估计式\eqref{eq:6}当中的系数，得到4组参数关于电流的二次函数。将估计结果代入式\eqref{eq:6}，计算各实测电流下拟合曲线与原始数据的$\mathrm{MRE}$，各电流下$\mathrm{MRE}$均保持在较低水平，表明模型在$20\,\text{A}$至$100\,\text{A}$区间具有可靠的预测精度。
\subsubsection{问题二55A放电曲线展示}
将 $I=55\,\text{A}$ 代入式\eqref{eq:7}，得到电流强度为$55\,\text{A}$时的放电曲线函数 $U=f(t,55)$。以时间步长$\Delta t=6.28\,\text{min}$生成电压序列，绘制曲线图并附部分数据表（完整见附录）。
\begin{table}[htbp]
\centering
\setlength{\tabcolsep}{16pt} % 增大列间距，加宽表格
\caption{部分55A预测放电电压‑时间数据表}
\label{tab:4}
\begin{tabular}{cc|cc}
\toprule
时间/min & 预测电压/V & 时间/min & 预测电压/V \\
\midrule
157.0951586 & 10.51125546 & 257.6360601 & 10.46064020 \\
163.3789649 & 10.50682367 & 263.9198664 & 10.45867286 \\
169.6627713 & 10.50258583 & 270.2036728 & 10.45681429 \\
175.9465776 & 10.49853663 & 276.4874791 & 10.45505919 \\
182.2303840 & 10.49467073 & 282.7712855 & 10.45340221 \\
188.5141903 & 10.49098282 & 289.0550918 & 10.45183804 \\
194.7979967 & 10.48746757 & 295.3388982 & 10.45036135 \\
& & 301.6227045 & 10.44896683 \\
\bottomrule
\end{tabular}
\end{table}
\begin{figure}[H]
\centering
\includegraphics[width=0.9\textwidth]{55A.png}
\caption{55A条件下预测放电曲线}
\label{fig:2}
\end{figure}
\section{模型分析检验}
\subsection{灵敏度分析}
\subsection{误差分析}
\section{模型的评价与推广}
\subsection{模型的优点}
\subsection{模型的不足}
\subsection{模型的推广}

\subsection{字体加粗与斜体}

如果想强调部分内容，可以使用加粗的手段来实现。加粗字体可以用 \verb|\textbf{加粗}| 来实现。例如： \textbf{这是加粗的字体。 This is bold fonts} 。

中文字体没有斜体设计，但是英文字体有。\textit{斜体 Italics}。

\section{参考文献与引用}

参考文献对于一篇正式的论文来说是必不可少的，在建模中重要的参考文献当然应该列出。\LaTeX{}在这方面的功能也是十分强大的，下面进介绍一个比较简单的参考文献制作方法。有兴趣的可以学习 \verb|bibtex| 或 \verb|biblatex| 的使用。

\LaTeX{}的入门书籍可以看《\LaTeX{}入门》\cite{liuhaiyang2013latex}。这是一个简单的引用，用 \verb|\cite{bibkey}| 来完成。要引用成功，当然要维护好 bibitem 了。下面是个简单的例子。



%参考文献
\begin{thebibliography}{9}%宽度9
    \bibitem[1]{liuhaiyang2013latex}
    刘海洋.
    \newblock \LaTeX {}入门\allowbreak[J].
    \newblock 电子工业出版社, 北京, 2013.
    \bibitem[2]{mathematical-modeling}
    全国大学生数学建模竞赛论文格式规范 (2020 年 8 月 25 日修改).
\end{thebibliography}

\newpage
%附录
\begin{appendices}

\section{模板所用的宏包}
\begin{table}[htbp]
    \centering
    \caption{宏包罗列}
    \begin{tabular}{ccccc}
        \toprule
        \multicolumn{5}{c}{模板中已经加载的宏包} \\
        \midrule
        amsbsy & amsfonts & {amsgen} & {amsmath} & {amsopn} \\
        amssymb & amstext & {appendix} & {array} & {atbegshi} \\
        atveryend & auxhook & {bigdelim} & {bigintcalc} & {bigstrut} \\
        bitset & bm    & {booktabs} & {calc} & {caption} \\
        caption3 & CJKfntef & {cprotect} & {ctex} & {ctexhook} \\
        ctexpatch & enumitem & {etexcmds} & {etoolbox} & {everysel} \\
        expl3 & fix-cm & {fontenc} & {fontspec} & {fontspec-xetex} \\
        geometry & gettitlestring & {graphics} & {graphicx} & {hobsub} \\
        hobsub-generic & hobsub-hyperref & {hopatch} & {hxetex} & {hycolor} \\
        hyperref & ifluatex & {ifpdf} & {ifthen} & {ifvtex} \\
        ifxetex & indentfirst & {infwarerr} & {intcalc} & {keyval} \\
        kvdefinekeys & kvoptions & {kvsetkeys} & {l3keys2e} & {letltxmacro} \\
        listings & longtable & {lstmisc} & {ltcaption} & {ltxcmds} \\
        multirow & nameref & {pdfescape} & {pdftexcmds} & {refcount} \\
        rerunfilecheck & stringenc & {suffix} & {titletoc} & {tocloft} \\
        trig  & ulem  & {uniquecounter} & {url} & {xcolor} \\
        xcolor-patch & xeCJK & {xeCJKfntef} & {xeCJK-listings} & {xparse} \\
        xtemplate & zhnumber &       &       &  \\
        \bottomrule
    \end{tabular}%
    \label{tab:addlabel}%
\end{table}%

以上宏包都已经加载过了，不要重复加载它们。

\section{排队算法--matlab 源程序}

\begin{lstlisting}[language=matlab]
kk=2;[mdd,ndd]=size(dd);
while ~isempty(V)
[tmpd,j]=min(W(i,V));tmpj=V(j);
for k=2:ndd
[tmp1,jj]=min(dd(1,k)+W(dd(2,k),V));
tmp2=V(jj);tt(k-1,:)=[tmp1,tmp2,jj];
end
tmp=[tmpd,tmpj,j;tt];[tmp3,tmp4]=min(tmp(:,1));
if tmp3==tmpd, ss(1:2,kk)=[i;tmp(tmp4,2)];
else,tmp5=find(ss(:,tmp4)~=0);tmp6=length(tmp5);
if dd(2,tmp4)==ss(tmp6,tmp4)
ss(1:tmp6+1,kk)=[ss(tmp5,tmp4);tmp(tmp4,2)];
else, ss(1:3,kk)=[i;dd(2,tmp4);tmp(tmp4,2)];
end;end
dd=[dd,[tmp3;tmp(tmp4,2)]];V(tmp(tmp4,3))=[];
[mdd,ndd]=size(dd);kk=kk+1;
end; S=ss; D=dd(1,:);
 \end{lstlisting}

 \section{规划解决程序--lingo源代码}

\begin{lstlisting}[language=c]
kk=2;
[mdd,ndd]=size(dd);
while ~isempty(V)
    [tmpd,j]=min(W(i,V));tmpj=V(j);
for k=2:ndd
    [tmp1,jj]=min(dd(1,k)+W(dd(2,k),V));
    tmp2=V(jj);tt(k-1,:)=[tmp1,tmp2,jj];
end
    tmp=[tmpd,tmpj,j;tt];[tmp3,tmp4]=min(tmp(:,1));
if tmp3==tmpd, ss(1:2,kk)=[i;tmp(tmp4,2)];
else,tmp5=find(ss(:,tmp4)~=0);tmp6=length(tmp5);
if dd(2,tmp4)==ss(tmp6,tmp4)
    ss(1:tmp6+1,kk)=[ss(tmp5,tmp4);tmp(tmp4,2)];
else, ss(1:3,kk)=[i;dd(2,tmp4);tmp(tmp4,2)];
end;
end
    dd=[dd,[tmp3;tmp(tmp4,2)]];V(tmp(tmp4,3))=[];
    [mdd,ndd]=size(dd);
    kk=kk+1;
end;
S=ss;
D=dd(1,:);
 \end{lstlisting}
\end{appendices}

\end{document} 